\section{Adaptive three-point kernels in the affine shear}
\label{sec:affine-adaptive-common-kernel}

The fixed-template argument of
Section~\ref{sec:affine-determinant-layer-entropy} does not allow the three
arm certificates to vary from point to point.  We now isolate the part of
the determinant mechanism that survives completely adaptive choices.

Retain the minimal-radical shear
\[
 U(h,k)=1+Rh,\qquad
 V(h,k)=1+R(h+ck),\qquad
 W(h,k)=1+R(h+bk),
\]
and let $p,q,r$ be three admissible parameters in $[1,M]^2$.  At each point
$s$ choose arbitrary divisors
\[
 d_U(s)\mid U(s),\qquad d_V(s)\mid V(s),\qquad d_W(s)\mid W(s).
\]
Define the triplewise common moduli
\[
 g_Z=\gcd(d_Z(p),d_Z(q),d_Z(r))\quad(Z=U,V,W),\qquad
 G=g_Ug_Vg_W.
\]

\begin{lemma}[Sharp square-triangle determinant]
\label{lem:sharp-square-triangle}
If three points lie in an axis-parallel square of side $N$, then
\[
                 |\det(q-p,r-p)|\le N^2.
\]
The constant is sharp.
\end{lemma}

\begin{proof}
Translate the square to $[0,N]^2$ and order the three first coordinates as
$x_1,x_1+s,x_1+s+t$, where $s,t\ge0$ and $s+t\le N$.  If the corresponding
second coordinates are $Y_1,Y_2,Y_3$, translation gives
\[
 D=tY_1-(s+t)Y_2+sY_3.
\]
The positive contribution and the absolute value of the negative
contribution are each at most $(s+t)N\le N^2$.  A permutation changes only
the sign.  Three corners attain equality.
\end{proof}

At any one admissible point the three arms are pairwise coprime and each is
coprime to $R$.  For a difference vector $(x,y)$ between two of the points,
cancellation of $R$ therefore gives
\[
 g_U\mid x,\qquad g_V\mid x+cy,\qquad g_W\mid x+by.
\]
Applying this to both $q-p$ and $r-p$ and using
\[
 \det((x,y),(x',y'))
 =(x+cy)y'-y(x'+cy')
 =(x+by)y'-y(x'+by')
\]
shows that every $g_Z$ divides the determinant.  The three moduli are
pairwise coprime, hence their product divides it.

\begin{theorem}[Adaptive common-kernel determinant bound]
\label{thm:adaptive-common-kernel-box}
If $p,q,r\in[1,M]^2$ are noncollinear, then
\[
                            G\le(M-1)^2.
\]
\end{theorem}

\begin{proof}
The preceding congruences give
$G\mid\det(q-p,r-p)$.  Noncollinearity makes the determinant nonzero, so
$G$ is at most its absolute value.  Apply
Lemma~\ref{lem:sharp-square-triangle} to the square of side $M-1$.
\end{proof}

This theorem allows unrelated certificates at all three points.  What it
does not provide is a density theorem forcing a triple whose three
triplewise gcds have large product.  That is now the precise positive
selection problem.

There are also useful canonical constant improvements.  Put
\[
 M=\left\lfloor\frac{c^6}{4R}\right\rfloor,\qquad c\ge9,
 \qquad E(n)=\frac n{\rad(n)}.
\]
For every $(h,k)\in[1,M]^2$,
\begin{equation}\label{eq:canonical-factor-three-arms}
                 3U\le c^6,\qquad 3V\le c^7,\qquad 3W\le c^7.
\end{equation}
Indeed, $4RM\le c^6$.  The short-arm estimate follows from
$12+12Rh\le12+3c^6\le4c^6$.  For a long arm,
\[
 4(3V)\le12+3c^6(c+1)\le4c^7,
\]
where the last inequality is $12\le(c-3)c^6$; $b\le c$ gives the same
bound for $W$.

Suppose the affine point meets the necessary exceptional threshold
\begin{equation}\label{eq:adaptive-excess-threshold}
              Rc^{14}<8192E(U)E(V)E(W).
\end{equation}
Since $3E(U)\le c^6$, cancellation gives
\[
                   3Rc^8<8192E(V)E(W).
\]
Thus at least one long arm $Z\in\{V,W\}$ satisfies
\begin{equation}\label{eq:long-arm-square-excess}
                   3Rc^8<8192E(Z)^2.
\end{equation}
Combining $Z=\rad(Z)E(Z)$, $3Z\le c^7$, and
\eqref{eq:long-arm-square-excess} yields
\begin{equation}\label{eq:long-arm-radical-square-support}
             27R\rad(Z)^2<8192c^6.
\end{equation}

For the full repeated kernel
\[
 K(n)=\prod_{p^e\parallel n,\ e\ge2}p^e,\qquad
 K=K(U)K(V)K(W),
\]
pairwise coprimality gives
\[
 K=\rad(K)E(U)E(V)E(W).
\]
Moreover, \eqref{eq:canonical-factor-three-arms} gives
$27K\le27UVW\le c^{20}$.  Multiplying
\eqref{eq:adaptive-excess-threshold} by $27\rad(K)$ and cancelling
$c^{14}$ proves the improved support bound
\begin{equation}\label{eq:full-kernel-factor-27-support}
                 27R\rad(K)<8192c^6.
\end{equation}

The geometric alternatives have been tested on actual parameters.  For
the primitive seed $(a,b,c)=(1,8,9)$ one has $R=6$ and $M=22143$.
The admissible points $(20,1)$ and $(20,2)$ have arm triples
\[
 (121,175,169),\qquad(121,229,217).
\]
They share $K(U)=121$, have distance one and vertical difference, while
\[
       121>(b+1)(c+1)=90,\qquad121>(c+1)^2=100.
\]
Taking divisor certificates $(121,25,169)$ and $(121,1,1)$ gives common
gcds $(121,1,1)$ and product $121$.  Thus this is a full-premise in-box
counterexample to the precise local implication that deletes the
zero-first-coordinate branch and concludes
$G\le(c+1)^2\lVert p-q\rVert_\infty^3=100$.

Likewise $(160,1),(160,2),(160,3)$ are admissible, collinear and share
$U=K(U)=961$.  Their diameter is two and
\[
                    961>(c+1)^2 2^3=800.
\]
Choosing certificates $(961,1,1)$ at every point gives triplewise gcds
$(961,1,1)$ and satisfies all six difference congruences.  Hence the precise
full-premise local cubic implication to noncollinearity is false.  This is
not a counterexample to Theorem~\ref{thm:adaptive-common-kernel-box}, nor to
its numerical box-wide conclusion: $961<(22143-1)^2$.  It refutes only the
local diameter-scale strengthening and the invalid step from divisibility
to a size bound when the determinant may be zero.

Lean verifies the sharp determinant lemma for every point order, the
derivation of the adaptive congruence interface from actual pointwise arm
divisors and triplewise gcds, all three bounds in
\eqref{eq:canonical-factor-three-arms}, the long-arm dichotomy,
\eqref{eq:long-arm-radical-square-support},
the arithmetic seam of \eqref{eq:full-kernel-factor-27-support}, and the
negations of both exact local implications.  No adaptive common-kernel density
input is assumed.  The affine route therefore remains active at the
triple-selection gate.
